\section{Future Work}
\label{sec:future-work}

This section outlines potential enhancements and research directions to extend the capabilities, performance, and applicability of the distributed MPI WiFi implementation.

\subsection{Performance Optimizations}

\subsubsection{Non-Blocking MPI Communication}

\paragraph{Motivation}:
Current synchronous MPI operations introduce significant overhead (36\% of total simulation time). Non-blocking communication could overlap computation with network transfer.

\paragraph{Proposed Implementation}:
\begin{lstlisting}[style=cpp]
// Current blocking approach
void RemoteYansWifiChannelStub::Send(...) {
    MpiInterface::SendMsg(&msg, size, 0, TX_REQUEST_TAG);  // Blocks
}

// Non-blocking alternative
void RemoteYansWifiChannelStub::Send(...) {
    MPI_Request request;
    MPI_Isend(&msg, size, MPI_BYTE, 0, TX_REQUEST_TAG, 
              MPI_COMM_WORLD, &request);
    m_pendingRequests.push_back(request);  // Track for completion
}

void RemoteYansWifiChannelStub::ProcessPendingRequests() {
    // Poll for completed sends
    for (auto& req : m_pendingRequests) {
        int flag;
        MPI_Test(&req, &flag, MPI_STATUS_IGNORE);
        if (flag) {
            // Send completed, remove from pending
        }
    }
}
\end{lstlisting}

\paragraph{Expected Benefits}:
\begin{itemize}
    \item Reduce MPI overhead from 36\% to ~15-20\%
    \item Improve real-time speedup from 0.78× to ~0.90×
    \item Enable larger simulations by reducing bottleneck
\end{itemize}

\paragraph{Challenges}:
\begin{itemize}
    \item Request tracking and completion polling
    \item Memory management for in-flight messages
    \item Ensuring causality preservation
\end{itemize}

\paragraph{Estimated Effort}: 3-4 weeks development + 2 weeks testing

\subsubsection{Message Batching}

\paragraph{Motivation}:
Sending individual packets as separate MPI messages is inefficient. Batching multiple transmissions into single messages reduces MPI call overhead.

\paragraph{Design Approach}:
\begin{lstlisting}[style=cpp]
class MessageBatcher {
public:
    void AddTransmission(uint32_t deviceId, Ptr<WifiPpdu> ppdu, ...);
    void Flush();  // Send batch to rank 0
    
private:
    struct BatchEntry {
        uint32_t deviceId;
        SerializedPpdu ppdu;
        MobilityInfo mobility;
        double txPowerDbm;
    };
    
    std::vector<BatchEntry> m_batch;
    Time m_batchInterval{MicroSeconds(100)};  // Accumulate for 100μs
};
\end{lstlisting}

\paragraph{Batching Strategy}:
\begin{enumerate}
    \item Accumulate transmissions for configurable interval (e.g., 100 μs)
    \item Send batch as single MPI message
    \item Rank 0 processes all transmissions in batch
\end{enumerate}

\paragraph{Expected Impact}:
\begin{table}[htbp]
\centering
\caption{Message batching projected performance}
\begin{tabular}{@{}lccc@{}}
\toprule
\textbf{Batch Size} & \textbf{MPI Messages} & \textbf{Overhead} & \textbf{Speedup} \\
\midrule
1 (current) & 944K & 36\% & 0.78× \\
10 & 94K & 22\% & 0.86× \\
50 & 19K & 12\% & 0.93× \\
100 & 9K & 8\% & 0.96× \\
\bottomrule
\end{tabular}
\end{table}

\paragraph{Trade-off}:
Larger batches increase latency (devices wait for batch to fill). Optimal batch size balances throughput and latency.

\paragraph{Estimated Effort}: 2-3 weeks

\subsubsection{Compression for Large Messages}

\paragraph{Motivation}:
Serialized PPDU messages can be large (>10 KB for aggregated frames). Compressing payloads reduces network transfer time.

\paragraph{Implementation Sketch}:
\begin{lstlisting}[style=cpp]
#include <zlib.h>

void CompressAndSend(const MpiMessage& msg) {
    std::vector<uint8_t> compressed;
    CompressData(msg.payload, compressed);
    
    MpiMessage compressedMsg = msg;
    compressedMsg.payload = compressed;
    compressedMsg.isCompressed = true;
    
    MpiInterface::SendMsg(&compressedMsg, ...);
}
\end{lstlisting}

\paragraph{Expected Compression Ratios}:
\begin{itemize}
    \item WiFi headers: 5:1 (repetitive structure)
    \item IP/UDP payload: 2:1 to 10:1 (depends on data)
    \item Overall: 3-4× size reduction
\end{itemize}

\paragraph{Trade-off}:
Compression adds CPU overhead (~5-10\%) but reduces network transfer time (~60-70\%). Net benefit depends on network vs. CPU bottleneck.

\paragraph{Estimated Effort}: 1-2 weeks

\subsection{Scalability Enhancements}

\subsubsection{Multi-Channel Processor Architecture}

\paragraph{Problem}:
Single channel processor (rank 0) becomes bottleneck beyond ~200 devices.

\paragraph{Solution - Hierarchical Channel Design}:

\begin{figure}[htbp]
\centering
\begin{verbatim}
                     Master Channel (Rank 0)
                    /          |          \
                   /           |           \
          Sub-channel 1  Sub-channel 2  Sub-channel 3
            (Rank 1)       (Rank 2)       (Rank 3)
              |                |              |
           Devices          Devices        Devices
         (Ranks 4-7)      (Ranks 8-11)   (Ranks 12-15)
\end{verbatim}
\caption{Hierarchical multi-channel processor architecture}
\end{figure}

\paragraph{Design Principles}:
\begin{enumerate}
    \item \textbf{Spatial partitioning}: Divide area into geographic zones
    \item \textbf{Local propagation}: Sub-channels handle intra-zone transmissions
    \item \textbf{Cross-zone coordination}: Master channel handles inter-zone propagation
    \item \textbf{Dynamic load balancing}: Redistribute zones based on density
\end{enumerate}

\paragraph{Implementation Approach}:
\begin{lstlisting}[style=cpp]
class SubChannelProcessor {
public:
    void ProcessLocalTransmission(uint32_t deviceId, ...);
    void ProcessCrossZone(uint32_t deviceId, ...);  // Forward to master
    
private:
    Box m_zone;  // Geographic boundary
    std::set<uint32_t> m_localDevices;
};

class MasterChannelProcessor {
public:
    void DistributeTransmission(uint32_t deviceId, ...);
    
private:
    std::map<uint32_t, uint32_t> m_deviceToSubChannel;
};
\end{lstlisting}

\paragraph{Expected Scalability Improvement}:
\begin{table}[htbp]
\centering
\caption{Hierarchical vs. centralized channel scalability}
\begin{tabular}{@{}lccc@{}}
\toprule
\textbf{Devices} & \textbf{Centralized} & \textbf{Hierarchical (4 sub)} & \textbf{Improvement} \\
\midrule
100 & 0.55× & 0.68× & 1.24× \\
200 & 0.42× & 0.59× & 1.40× \\
500 & 0.28× & 0.48× & 1.71× \\
1000 & N/A (too slow) & 0.35× & Enables new scale \\
\bottomrule
\end{tabular}
\end{table}

\paragraph{Challenges}:
\begin{itemize}
    \item Zone boundary handling: Devices near boundaries interact with multiple zones
    \item Load balancing: Uneven device distribution requires dynamic zone resizing
    \item Mobility: Devices moving between zones trigger reassignment
\end{itemize}

\paragraph{Estimated Effort}: 6-8 weeks (significant architectural change)

\subsubsection{Spatial Indexing for Range-Limited Propagation}

\paragraph{Motivation}:
WiFi has finite range (~100-300m). Computing propagation to all devices wastes computation for distant pairs.

\paragraph{Approach}:
Use spatial index (e.g., R-tree, grid) to query only nearby devices.

\begin{lstlisting}[style=cpp]
class SpatialIndex {
public:
    void Insert(uint32_t deviceId, Vector position);
    std::vector<uint32_t> QueryWithinRange(Vector position, double range);
};

void WifiChannelMpiProcessor::ProcessTransmission(...) {
    Vector txPosition = GetPosition(txDeviceId);
    
    // Only compute for devices within max WiFi range
    auto nearbyDevices = m_spatialIndex.QueryWithinRange(
        txPosition, 300.0  // meters
    );
    
    for (auto rxId : nearbyDevices) {
        double rxPower = CalculatePropagation(txPosition, GetPosition(rxId), ...);
        SendRxNotification(rxId, rxPower, ...);
    }
}
\end{lstlisting}

\paragraph{Complexity Reduction}:
\begin{align}
\text{Centralized} &: O(N^2) \text{ propagation calculations} \notag \\
\text{Range-limited} &: O(N \times k) \text{ where } k \ll N \notag
\end{align}

For V2X scenario with 200 vehicles, if only ~20 vehicles within range:
\begin{align}
\text{Calculations} &: 200 \times 199 = 39,800 \rightarrow 200 \times 20 = 4,000 \notag \\
\text{Speedup} &: 9.95\times
\end{align}

\paragraph{Estimated Effort}: 3-4 weeks

\subsubsection{GPU-Accelerated Propagation Calculation}

\paragraph{Motivation}:
Propagation calculations are embarrassingly parallel—each device pair independent.

\paragraph{Approach}:
Offload propagation to GPU using CUDA or OpenCL.

\begin{lstlisting}[style=cpp]
// CPU version (current)
for (auto rxId : devices) {
    double rxPower = m_propagationModel->CalcRxPower(
        txPower, txMobility, rxMobility
    );
}

// GPU version (proposed)
__global__ void CalculatePropagationKernel(
    DeviceInfo* devices, int numDevices,
    TransmissionInfo tx, double* results
) {
    int rxId = blockIdx.x * blockDim.x + threadIdx.x;
    if (rxId < numDevices) {
        results[rxId] = LogDistanceLoss(
            tx.power, tx.position, devices[rxId].position, ...
        );
    }
}
\end{lstlisting}

\paragraph{Expected Speedup}:
\begin{itemize}
    \item CPU: Sequential, ~1M calculations/second
    \item GPU (NVIDIA T4): Parallel, ~100M calculations/second
    \item \textbf{Speedup}: 100× for propagation calculations
\end{itemize}

\paragraph{Impact on Overall Performance}:
Propagation is 27\% of total time. 100× speedup on 27\% → 1.37× overall speedup.

\paragraph{Challenges}:
\begin{itemize}
    \item GPU memory management for device positions
    \item Host-device data transfer overhead
    \item Requires CUDA-capable hardware
\end{itemize}

\paragraph{Estimated Effort}: 4-5 weeks

\subsection{Feature Extensions}

\subsubsection{Dynamic Mobility Support}

\paragraph{Current Limitation}:
Mobility models must be pre-defined and synchronized across ranks.

\paragraph{Proposed Solution}:
Centralized mobility manager broadcasts position updates.

\begin{lstlisting}[style=cpp]
class MobilityManager {
public:
    void UpdatePosition(uint32_t deviceId, Vector newPosition);
    void BroadcastPositionUpdates();
    
private:
    std::map<uint32_t, Vector> m_positions;
    Time m_updateInterval{MilliSeconds(100)};  // Broadcast every 100ms
};
\end{lstlisting}

\paragraph{Message Flow}:
\begin{enumerate}
    \item Mobility models update positions locally
    \item Device ranks send position updates to rank 0
    \item Rank 0 broadcasts consolidated updates to all ranks
    \item All ranks maintain synchronized position view
\end{enumerate}

\paragraph{Optimization}:
Only broadcast positions that changed significantly (> 1 meter).

\paragraph{Estimated Effort}: 3 weeks

\subsubsection{Heterogeneous Propagation Models}

\paragraph{Motivation}:
Mixed environments (e.g., indoor + outdoor) require different models.

\paragraph{Design}:
Per-pair propagation model selection based on position.

\begin{lstlisting}[style=cpp]
class HybridPropagationModel : public PropagationLossModel {
public:
    void SetIndoorRegion(Box region, Ptr<PropagationLossModel> model);
    void SetOutdoorModel(Ptr<PropagationLossModel> model);
    
    double CalcRxPower(double txPower, 
                       Ptr<MobilityModel> a, 
                       Ptr<MobilityModel> b) override {
        if (BothIndoor(a, b)) {
            return m_indoorModel->CalcRxPower(txPower, a, b);
        } else if (BothOutdoor(a, b)) {
            return m_outdoorModel->CalcRxPower(txPower, a, b);
        } else {
            return m_indoorOutdoorModel->CalcRxPower(txPower, a, b);
        }
    }
};
\end{lstlisting}

\paragraph{Estimated Effort}: 2-3 weeks

\subsubsection{Advanced Traffic Patterns}

\paragraph{Current Support}:
OnOff and UDP constant-rate traffic.

\paragraph{Proposed Additions}:
\begin{enumerate}
    \item \textbf{Video streaming}: Variable bitrate with I/P/B frames
    \item \textbf{Voice over WiFi}: VoIP codecs (G.711, Opus) with jitter requirements
    \item \textbf{Web traffic}: HTTP request-response patterns with think time
    \item \textbf{File transfer}: TCP with congestion control
\end{enumerate}

\paragraph{Implementation Example - Video Streaming}:
\begin{lstlisting}[style=cpp]
class VideoStreamApplication : public Application {
public:
    void SetFrameRate(double fps);
    void SetResolution(uint32_t width, uint32_t height);
    void SetCodec(VideoCodec codec);  // H.264, H.265, VP9
    
private:
    void SendIFrame();   // Intra-frame (large)
    void SendPFrame();   // Predicted frame (medium)
    void SendBFrame();   // Bidirectional frame (small)
};
\end{lstlisting}

\paragraph{Estimated Effort}: 2-3 weeks per application type

\subsubsection{Real-Time Monitoring Dashboard}

\paragraph{Motivation}:
Long simulations benefit from live monitoring to detect issues early.

\paragraph{Proposed Design}:
Web-based dashboard showing real-time metrics.

\begin{lstlisting}[style=cpp]
class MonitoringServer {
public:
    void ReportMetric(std::string name, double value);
    void StartWebServer(uint16_t port);
    
private:
    void HandleMetricsRequest(HttpRequest& req, HttpResponse& res);
    std::map<std::string, std::deque<TimeSeriesPoint>> m_timeSeries;
};

// Usage
MonitoringServer monitor;
monitor.StartWebServer(8080);

Simulator::Schedule(Seconds(1.0), [&]() {
    monitor.ReportMetric("rank0_cpu", GetCpuUsage(0));
    monitor.ReportMetric("mpi_queue_size", GetQueueSize());
});
\end{lstlisting}

\paragraph{Dashboard Features}:
\begin{itemize}
    \item CPU utilization per rank (line charts)
    \item Memory usage trends
    \item MPI message rates and latency
    \item Simulation progress (percentage complete)
    \item Live packet transmission visualization
\end{itemize}

\paragraph{Technologies}:
\begin{itemize}
    \item Backend: Simple HTTP server (microhttpd)
    \item Frontend: React + Chart.js
    \item Data format: JSON REST API
\end{itemize}

\paragraph{Estimated Effort}: 4-5 weeks

\subsection{Usability Improvements}

\subsubsection{Automated MPI Configuration}

\paragraph{Current Pain Point}:
Users must manually create hostfiles, configure SSH, set environment variables.

\paragraph{Proposed Solution}:
Setup wizard script.

\begin{lstlisting}[style=bash]
#!/bin/bash
# setup-mpi-cluster.sh

echo "=== NS-3 MPI WiFi Setup Wizard ==="

# Detect available machines
echo "Enter hostnames/IPs (one per line, Ctrl-D when done):"
HOSTS=()
while read -r host; do
    HOSTS+=("$host")
done

# Test SSH connectivity
for host in "${HOSTS[@]}"; do
    echo "Testing SSH to $host..."
    ssh -o ConnectTimeout=5 "$host" "echo OK" || {
        echo "❌ SSH failed to $host"
        exit 1
    }
done

# Generate hostfile
cat > ~/hosts.txt << EOF
$(for host in "${HOSTS[@]}"; do echo "$host slots=4"; done)
EOF

# Test MPI
echo "Testing MPI across cluster..."
mpirun --hostfile ~/hosts.txt -np ${#HOSTS[@]} hostname

echo "✅ MPI cluster ready!"
\end{lstlisting}

\paragraph{Estimated Effort}: 1-2 weeks

\subsubsection{Cloud Deployment Templates}

\paragraph{Motivation}:
Simplify AWS/GCP/Azure deployment.

\paragraph{Deliverables}:
\begin{enumerate}
    \item \textbf{Terraform scripts}: Infrastructure as code for cluster provisioning
    \item \textbf{Docker containers}: Pre-configured NS-3 + MPI environment
    \item \textbf{Kubernetes manifests}: Orchestrated multi-pod simulation
\end{enumerate}

\paragraph{Example - AWS Terraform}:
\begin{lstlisting}[language=json, caption={terraform/main.tf}]
resource "aws_instance" "ns3_node" {
  count         = var.num_instances
  ami           = "ami-0c55b159cbfafe1f0"  # Ubuntu 22.04
  instance_type = "t3.xlarge"
  
  tags = {
    Name = "ns3-mpi-node-${count.index}"
  }
  
  user_data = file("setup-ns3.sh")  # Installs NS-3, MPI
}
\end{lstlisting}

\paragraph{Estimated Effort}: 3-4 weeks

\subsection{Research Directions}

\subsubsection{Optimistic Synchronization}

\paragraph{Concept}:
Instead of synchronous MPI (conservative), allow ranks to proceed optimistically and rollback if causality violations detected.

\paragraph{Potential Benefits}:
\begin{itemize}
    \item Eliminate blocking MPI overhead
    \item Better CPU utilization
    \item Potential 2-3× speedup
\end{itemize}

\paragraph{Challenges}:
\begin{itemize}
    \item State checkpointing for rollback
    \item Causality violation detection
    \item Complex debugging when rollbacks occur
\end{itemize}

\paragraph{Research Questions}:
\begin{itemize}
    \item What is optimal lookahead time for WiFi?
    \item How often do rollbacks occur in typical scenarios?
    \item Does benefit justify complexity?
\end{itemize}

\paragraph{Estimated Effort}: 3-4 months (research project)

\subsubsection{Machine Learning for Load Balancing}

\paragraph{Concept}:
Use ML to predict optimal device-to-rank assignments based on mobility patterns and traffic.

\paragraph{Approach}:
\begin{enumerate}
    \item Collect training data from various scenarios
    \item Train model: Input (device positions, traffic) → Output (rank assignments)
    \item Deploy model to dynamically rebalance during simulation
\end{enumerate}

\paragraph{Potential Impact}:
10-20\% performance improvement by preventing hotspots.

\paragraph{Estimated Effort}: 4-6 months (research + implementation)

\subsection{Integration and Deployment}

\subsubsection{NS-3 Mainline Integration}

\paragraph{Goal}:
Contribute implementation to official NS-3 repository.

\paragraph{Requirements}:
\begin{enumerate}
    \item Code review and cleanup
    \item Comprehensive unit tests
    \item Documentation (Doxygen, tutorial)
    \item Backward compatibility (MPI optional, not required)
    \item CI/CD integration
\end{enumerate}

\paragraph{Estimated Effort}: 2-3 months

\subsubsection{Open-Source Release}

\paragraph{Deliverables}:
\begin{itemize}
    \item GitHub repository with examples
    \item Docker images for easy setup
    \item YouTube tutorial videos
    \item Discourse forum for community support
\end{itemize}

\paragraph{Estimated Effort}: 1-2 months

\subsection{Priority Roadmap}

Based on impact and effort, we recommend the following priority order:

\begin{table}[htbp]
\centering
\caption{Future work priority roadmap}
\label{tab:future-work-roadmap}
\small
\begin{tabular}{@{}clcc@{}}
\toprule
\textbf{Priority} & \textbf{Enhancement} & \textbf{Impact} & \textbf{Effort} \\
\midrule
1 (High) & Message batching & High & 2-3 weeks \\
2 (High) & Spatial indexing & High & 3-4 weeks \\
3 (High) & Automated setup wizard & Medium & 1-2 weeks \\
4 (Medium) & Non-blocking MPI & High & 3-4 weeks \\
5 (Medium) & Dynamic mobility & Medium & 3 weeks \\
6 (Medium) & Cloud templates & Medium & 3-4 weeks \\
7 (Low) & Multi-channel processor & Very High & 6-8 weeks \\
8 (Low) & GPU propagation & High & 4-5 weeks \\
9 (Low) & Real-time dashboard & Medium & 4-5 weeks \\
10 (Research) & Optimistic sync & Very High & 3-4 months \\
\bottomrule
\end{tabular}
\end{table}

\textbf{Recommended 6-month plan}:
\begin{enumerate}
    \item \textbf{Months 1-2}: Message batching, spatial indexing, setup wizard (quick wins)
    \item \textbf{Months 3-4}: Non-blocking MPI, dynamic mobility, cloud templates (core improvements)
    \item \textbf{Months 5-6}: Multi-channel processor or GPU propagation (scalability breakthrough)
\end{enumerate}

This roadmap would deliver immediate performance gains (months 1-2), solidify the implementation (months 3-4), and enable next-level scalability (months 5-6).
